\section*{Balance solenoidal, acarreo y terminal enriquecido}

Este módulo refuerza el \texttt{OWNER\_RECTOR} Navier--Stokes: conserva en
una sola cadena el acarreo de Galerkin, el balance causal, el Gram terminal,
Stein, la telescopía, la inyección raíz, el retorno y el paso global. La carta
KYP aparece sólo después y queda fuera de esa cadena.

Para los cortes de Galerkin \(m\leq n\), el acarreo conservado es
\begin{equation}\label{eq:amp-ns-owner-acarreo-galerkin}
 C_{m\leftarrow n}
 :=P_mB(u_n,u_n)-B_m(P_m u_n,P_m u_n).
\end{equation}
La lectura nonádica de este estado está tipada por
\begin{equation}\label{eq:amp-ns-owner-e9j9}
 E_9J_9=I,\qquad P_9=J_9E_9,\qquad Q_{\rm micro}=I-P_9,
\end{equation}
y conserva el balance causal
\begin{equation}\label{eq:amp-ns-owner-balance-causal}
 \boxed{
 \mathbb E_9\|\Xi_{m+1}\|^2+\mathbb E_9\ell_m
 =\|\Xi_m\|^2,\qquad \ell_m\geq0.}
\end{equation}
El campo, la proyección solenoidal, el acarreo, la frontera, la ruta y las
pérdidas positivas permanecen así en un mismo estado antes de construir el
almacenamiento.

\section*{Gram truncado del estado completo e inyección raíz uniforme}

La realización que gobierna el balance no comienza por la desigualdad KYP.
Sean \(\Gamma_{M,j}^{\rm term}\) la transición de un paso del estado
PC--Fock enriquecido, \(\Gamma_{M,j:k}^{\rm term}\) su composición entre
las profundidades \(j\) y \(k\), y
\(L_{M,k,T}^{\rm term}\) la columna de pérdidas de grado uno, grado dos,
acarreo, innovación, archivo y memoria en el horizonte \([0,T]\). La
notación \(\mathbb E_{j:k}\) designa la suma ortogonal con los pesos de
las ramas del refinamiento entre esos niveles. Para \(j\leq J\) se define
la columna finita
\begin{equation}\label{eq:amp-ns-owner-columna-gram-truncado}
 \begin{aligned}
 \mathbf G_{M,j;J,T}^{\rm term}x
 :={}&
 \bigl(\mathcal E_{M,J,T}^{\rm term}
       \Gamma_{M,j:J}^{\rm term}x\bigr)_{\mathbb E_{j:J}}\\
 &\oplus
 \bigoplus_{k=j}^{J-1}
 \bigl(L_{M,k,T}^{\rm term}
       \Gamma_{M,j:k}^{\rm term}x\bigr)_{\mathbb E_{j:k}},
 \qquad
 U_{M,j;J,T}^{\rm term}
 :=(\mathbf G_{M,j;J,T}^{\rm term})^*
    \mathbf G_{M,j;J,T}^{\rm term}.
 \end{aligned}
\end{equation}
En particular, el almacenamiento es un Gram truncado construido con el
terminal y las pérdidas futuras del mismo estado. Su forma cuadrática es
\begin{equation}\label{eq:amp-ns-owner-gram-truncado}
 \begin{aligned}
 \langle x,U_{M,j;J,T}^{\rm term}x\rangle
 ={}&\mathbb E_{j:J}
 \|\mathcal E_{M,J,T}^{\rm term}
       \Gamma_{M,j:J}^{\rm term}x\|^2\\
 &+\sum_{k=j}^{J-1}\mathbb E_{j:k}
 \|L_{M,k,T}^{\rm term}
       \Gamma_{M,j:k}^{\rm term}x\|^2 .
 \end{aligned}
\end{equation}

\begin{lemma}[Identidad de Stein del Gram truncado]
\label{lem:amp-ns-owner-stein-truncado}
Para todo \(j<J\),
\begin{equation}\label{eq:amp-ns-owner-stein-truncado}
 U_{M,j;J,T}^{\rm term}
 =(\Gamma_{M,j}^{\rm term})^*
   U_{M,j+1;J,T}^{\rm term}\Gamma_{M,j}^{\rm term}
  +(L_{M,j,T}^{\rm term})^*L_{M,j,T}^{\rm term}.
\end{equation}
Por iteración se obtiene, conservando el sumando terminal,
\begin{equation}\label{eq:amp-ns-owner-telescopio-truncado}
 \begin{aligned}
 U_{M,0;J,T}^{\rm term}
 ={}&(\Gamma_{M,0:J}^{\rm term})^*
      U_{M,J;J,T}^{\rm term}\Gamma_{M,0:J}^{\rm term}\\
 &+\sum_{j=0}^{J-1}
   (\Gamma_{M,0:j}^{\rm term})^*
   (L_{M,j,T}^{\rm term})^*L_{M,j,T}^{\rm term}
   \Gamma_{M,0:j}^{\rm term}.
 \end{aligned}
\end{equation}
\end{lemma}

\begin{proof}
En \eqref{eq:amp-ns-owner-columna-gram-truncado}, se separa el sumando
\(k=j\). Los restantes sumandos, incluido el terminal, son exactamente la
columna de nivel \(j+1\) precedida por
\(\Gamma_{M,j}^{\rm term}\). La ortogonalidad de la suma y la propiedad de
torre de \(\mathbb E_{j:k}\) dan
\eqref{eq:amp-ns-owner-stein-truncado}. La repetición del mismo
desdoblamiento da \eqref{eq:amp-ns-owner-telescopio-truncado}; no se
descarta
\((\Gamma_{M,0:J}^{\rm term})^*U_{M,J;J,T}^{\rm term}
\Gamma_{M,0:J}^{\rm term}\).
\end{proof}

\begin{theorem}[Inyección raíz y presupuesto uniforme]
\label{thm:amp-ns-inyeccion-raiz-uniforme}
La raíz del realizador se transporta a cada corte y profundidad mediante
una inclusión isométrica:
\begin{equation}\label{eq:amp-ns-inyeccion-raiz-comun}
 \begin{aligned}
 x_{M,0}^{\rm data}
 &:=\mathcal V_M^{\rm full}(P_Mu_0)\oplus g_M^F,\\
 \mathbf G_{M,0;J,T}^{\rm term}x_{M,0}^{\rm data}
 &=\mathcal J_{M,J,T}^{\rm term}(u_0,f),\\
 \mathcal J_{M,J,T}^{\rm term}
 &:=I_{0\to M,J}\mathcal J_T^{\rm term},
 & I_{0\to M,J}^*I_{0\to M,J}&=I.
 \end{aligned}
\end{equation}
Para los datos de \eqref{eq:realizaciones-ns-dominio}, existe
\(C_{{\rm HMT},T}\), independiente de \(M\) y \(J\), tal que
\begin{equation}\label{eq:amp-ns-inyeccion-raiz-cota}
 \sup_{M,J}
 \|\mathcal J_{M,J,T}^{\rm term}(u_0,f)\|^2
 \leq C_{{\rm HMT},T}
 +(1+2\widetilde M_{s,\beta,\nu}^{\,2})
  e^{R_{\rm F}^{\,2}\|u_0\|_{H^s}^2}
 +C_{\nu,s}\int_0^T\|f(t)\|_{H^{s-1}}^2\,dt .
\end{equation}
Equivalentemente, la cota es una cota del Gram terminal evaluado en
los datos:
\begin{equation}\label{eq:amp-ns-owner-gram-cota-uniforme}
 \sup_{M,J}
 \left\langle x_{M,0}^{\rm data},
 U_{M,0;J,T}^{\rm term}x_{M,0}^{\rm data}\right\rangle
 =
 \sup_{M,J}
 \|\mathcal J_{M,J,T}^{\rm term}(u_0,f)\|^2
 <\infty .
\end{equation}
La aplicación \(\mathcal J_T^{\rm term}\) depende sólo del estado inicial,
de sus dos hojas, del terminal heredado y del puerto de fuerza; no recibe
como argumento \(u_M(t)\) para \(0<t\leq T\).
\end{theorem}

\begin{proof}
La componente lineal y las memorias \(D/H\) se transportan
isométricamente. En la componente PC--Fock, la columna coherente es
\[
 \bigoplus_{n\geq0}
 \frac{R_{\rm F}^{\,n}}{\sqrt{n!}}\,u_0^{\odot n},
 \qquad
 \sum_{n\geq0}\frac{R_{\rm F}^{\,2n}}{n!}
       \|u_0\|_{H^s}^{2n}
 =e^{R_{\rm F}^{\,2}\|u_0\|_{H^s}^2}.
\]
Para justificar la uniformidad del lector cuadrático, se polariza
\begin{equation}\label{eq:amp-ns-lector-cuadratico-polarizado}
 \mathcal B_M(u,v):=
 \frac{1}{4\sqrt{\beta\nu}}\,
 A_M^{-1/2}\Lambda^sP_M\mathbb P
 \bigl((P_Mu\!\cdot\!\nabla)P_Mv
      +(P_Mv\!\cdot\!\nabla)P_Mu\bigr).
\end{equation}
En las bases transversales
\(e^s_{p,a}=\langle p\rangle^{-s}a\,e^{ip\cdot x}\), si
\(k=p+q\ne0\), sus filas de Fourier cumplen
\begin{equation}\label{eq:amp-ns-filas-fourier-uniformes}
 \|c^M_{k;p,q}\|
 \leq\frac{C}{\sqrt{\beta\nu}}\,
 \frac{\langle k\rangle^s(|p|+|q|)}
 {|k|\langle p\rangle^s\langle q\rangle^s},
 \qquad
 \sup_{M,k\ne0}\sum_{p+q=k}\|c^M_{k;p,q}\|^2<\infty .
\end{equation}
La segunda cota se obtiene separando \(|p|\leq2|k|\) de
\(|p|>2|k|\); en la cola domina
\(\sum_p\langle p\rangle^{2-2s}<\infty\) porque \(s>5/2\).
Cauchy--Schwarz en cada fila y la ortogonalidad de los bloques con distinto
\(k\) prolongan \(\mathcal B_M\) a
\[
 \mathfrak C_M:H^s_\sigma\odot_2H^s_\sigma\longrightarrow L^2_\sigma,
 \qquad
 \sup_M\|\mathfrak C_M\|
 \leq\widetilde M_{s,\beta,\nu}.
\]
Así, la constante usada en la raíz controla la suma de todos los pares
Fourier y no sólo cada coeficiente por separado. El puerto afín verifica
\[
 Q_f(T)\leq C_{\nu,s}
 \int_0^T\|f(t)\|_{H^{s-1}}^2\,dt .
\]
Por suma ortogonal de esas componentes,
\begin{equation}\label{eq:amp-ns-raiz-cota-por-componentes}
 \|\Xi_{M,0}\|^2
 \leq C_{{\rm HMT},T}
 +(1+2\widetilde M_{s,\beta,\nu}^{\,2})
 e^{R_{\rm F}^{\,2}\|u_0\|_{H^s}^2}.
\end{equation}
Las inclusiones de \eqref{eq:amp-ns-inyeccion-raiz-comun} no cambian esta
norma. Al sumar la cota del puerto se obtiene exactamente
\eqref{eq:amp-ns-inyeccion-raiz-cota}. Como todos los términos se evalúan
en \((u_0,f)\), la estimación precede a la trayectoria y no presupone la
cota global que después se deduce de
\eqref{eq:amp-ns-owner-telescopio-truncado}.
\end{proof}

El Gram terminal \eqref{eq:amp-ns-owner-gram-truncado}, su identidad de
Stein, la telescopía con terminal y la inyección raíz uniforme forman la
construcción propietaria. Ésta es la rama que gobierna el cierre y coincide
con la realización íntegra restaurada en
\cref{sec:ns-owner-closure,thm:nsclose-proprietary-full}.

\section*{Separación del control KYP finito no gobernante}

La factorización KYP finita correcta del
\cref{thm:ns-import-affine-stein} se conserva como control alternativo de
una carta reducida y no constituye una premisa del owner
\(U^{\rm own}\). No se utiliza la mutación firmada que pretendía
incorporar \(C_{m,M}^*C_{m,M}\) a un término \(P\succeq0\) y
recuperarlo después con signo negativo: al entrar en \(P\), ese
sumando aparece necesariamente con signo positivo. Por tanto, esa
factorización no interviene en el balance, el presupuesto raíz, el retorno
físico ni la conclusión Navier--Stokes.

\section*{Derivada temporal, presión y pegado de horizontes}

El tramo siguiente es consecuencia de la cota uniforme de trayectoria
obtenida por el Gram terminal propietario, su telescopía y el retorno
físico exacto. Se
dispone de la cota conjunta uniforme en
\(L^\infty(0,T;H^s)\cap L^2(0,T;H^{s+1})\).
La ecuación proyectada se escribe
\begin{equation}\label{eq:amp-ns-proyectada}
 \partial_tu_m
 =\nu\Delta u_m-P_mB(u_m,u_m)+P_mf.
\end{equation}
Para \(s>5/2\), el producto de Sobolev y la cota conjunta en
\(L^\infty(0,T;H^s)\cap L^2(0,T;H^{s+1})\) dan
\begin{equation}\label{eq:amp-ns-dt}
 \sup_m\|\partial_tu_m\|_{L^2(0,T;H^{s-1})}
 \leq C_T(u_0,f,\Xi_0).
\end{equation}
La estimación \eqref{eq:amp-ns-dt} y la compacidad espacial producen
convergencia fuerte en \(L^2(0,T;H^s)\), suficiente para pasar al término
bilineal sin perder el carry de Galerkin.

Recuperada la velocidad, la presión normalizada de media cero queda fijada
por
\begin{equation}\label{eq:amp-ns-presion}
 -\Delta p
 =\partial_i\partial_j(u_i u_j)-\nabla\!\cdot f,
 \qquad \int_{\mathbb T^3}p(t,x)\,dx=0.
\end{equation}
La presión es así una salida de la misma solución solenoidal; no es una
variable libre usada para seleccionar la rama regular.

\begin{proposition}[Pegado global por unicidad]
\label{prop:amp-ns-pegado}
Sean \(u^{(T)}\) las soluciones obtenidas en \([0,T]\) por el límite de
Galerkin. Si \(T_1<T_2\), entonces
\(u^{(T_2)}|_{[0,T_1]}=u^{(T_1)}\). La familia define una única solución
global suave para \(t>0\).
\end{proposition}

\begin{proof}
Ambas restricciones poseen los mismos datos y pertenecen al dominio de la
estimación uniforme. La ecuación de su diferencia, el control bilineal en
\(H^{s-1}\) y Grönwall implican igualdad en \([0,T_1]\). Los intervalos
finitos forman un sistema dirigido; la igualdad en los solapamientos define
la solución global. La regularización parabólica del teorema se aplica en
cada intervalo \([\varepsilon,T]\), \(\varepsilon>0\).
\end{proof}

El control de refutación es preciso: una pérdida de uniformidad en
\eqref{eq:amp-ns-dt}, dos límites diferentes en un mismo horizonte o una
presión que no satisfaga \eqref{eq:amp-ns-presion} romperían la cadena antes
de la conclusión global.

\paragraph{Conclusión propietaria.}
El Gram terminal propietario, Stein, la telescopía con terminal, su ledger,
el retorno físico exacto y la inyección raíz producen una cota \(H^s\) uniforme para
\(s>5/2\). Por tanto se demuestra la existencia, unicidad y regularidad
global de Navier--Stokes tridimensional periódico en el dominio HMT
registrado. Las cartas KYP finitas permanecen como
\texttt{CONTROL\_NO\_GOBERNANTE}.
